\documentclass[11pt,a4paper,landscape]{article}

% Atlas de arquitectura do Opcoda.
%
% Este documento e' a forma LaTeX do mesmo atlas que esta em Markdown em
% docs/atlas/*.md. As figuras vem de docs/atlas/figuras/*.tex, geradas por
% tools/mmd2tikz.py a partir de docs/atlas/diagrams/*.mmd. O .mmd e' a fonte
% da verdade; este ficheiro e' a montagem.
%
% A pagina e' deitada porque os diagramas de fluxo sao largos: uma cadeia de
% camadas nao cabe numa folha deitada com o texto a dobrar a largura. O
% texto corrido vai em duas colunas, que e' a medida que se le.
%
% Para gerar: powershell -File tools/build-atlas.ps1

\usepackage[T1]{fontenc}
\usepackage[utf8]{inputenc}
\usepackage[brazil]{babel}
\usepackage{lmodern}
\usepackage{newtxtext}
\usepackage{geometry}
\geometry{top=1.8cm,bottom=1.8cm,left=2cm,right=2cm}
\usepackage{setspace}
\setstretch{1.1}
\usepackage{graphicx}
\usepackage{booktabs}
\usepackage{array}
\usepackage{longtable}
\usepackage{multicol}
\usepackage{url}
\usepackage{hyperref}
\usepackage{tikz}
\usetikzlibrary{positioning,arrows.meta,fit,calc,shapes.misc,backgrounds,shapes.multipart}

\pagestyle{plain}
\setlength{\parindent}{0.8cm}
\setlength{\parskip}{0pt}
\setlength{\emergencystretch}{3cm}
\setlength{\columnsep}{0.8cm}

% As figuras sao largas. O TikZ gerado parte a cadeia de camadas em filas
% para caber na largura, e esta macro so encolhe o que ainda nao caiba.
\newsavebox{\figbox}
\newcommand{\maxfig}[1]{%
  \sbox{\figbox}{\input{figuras/#1.tex}}%
  \ifdim\wd\figbox>\linewidth
    \resizebox{\linewidth}{!}{\usebox{\figbox}}%
  \else
    \usebox{\figbox}%
  \fi
}

\newcolumntype{L}[1]{>{\raggedright\arraybackslash}p{#1}}
% As tabelas somam 25,7 cm de texto. Com o \tabcolsep por omissao, a soma
% das colunas passa do alinhamento e o LaTeX reclama com um overfull.
\setlength{\tabcolsep}{4pt}

\begin{document}

\begin{center}
\MakeUppercase{Opcoda: atlas da arquitectura}
\end{center}

\begin{flushright}
Mateus Costa de Oliveira Santos\\
Rafael Augusto Garcez Elisei\\
Prof. Me. Warner Brezolin (orientador)
\end{flushright}

\section*{1 Mapa do repositorio}

\begin{multicols}{2}
O Opcoda e' um sintetizador granular que transforma ficheiros executaveis do
formato Portable Executable em som. A thread de audio nunca le disco: a
leitura, a analise e a conversao do material acontecem na thread de
interface, e a unica passagem entre as duas e' uma fila circular sem trava
que carrega ponteiro cru.

O repositorio tem duas camadas. \texttt{opcoda\_core} e' C++20 puro e nao
inclui JUCE em lado nenhum; \texttt{opcoda\_plugin} e' a unica camada que
depende do framework e da thread de audio do host. O nucleo existe para
poder ser testado sem uma janela e sem um host.
\end{multicols}

\begin{center}
\begin{tabular}{@{}L{6.0cm}L{19.0cm}@{}}
\toprule
Pasta & Papel \\
\midrule
\path{src/opcoda\_core/pe} & Leitura defensiva do PE, conversao de bytes, selecao de regiao \\
\path{src/opcoda\_core/entropy} & Entropia de Shannon por janela \\
\path{src/opcoda\_core/dsp} & Motor granular, janelas, filtros, limitador \\
\path{src/opcoda\_core/rt} & Fila sem trava, transporte, notas, guard de alocacao \\
\path{src/opcoda\_core/view} & Rastro dos graos, sem JUCE \\
\path{src/opcoda\_plugin/source} & Processador JUCE, editor, grelha de bytes \\
\bottomrule
\end{tabular}
\end{center}

\section*{2 Contexto}

O Opcoda entra em cena por tres vias: o utilizador larga um executavel, o
host entrega MIDI e parametros, e o host entrega o estado de um projecto
guardado. Do outro lado, o plugin devolve audio, devolve estado e escreve
ficheiros de preset.

\begin{center}
\maxfig{context}
\end{center}

\begin{center}
\begin{tabular}{@{}L{3.6cm}L{6.6cm}L{6.4cm}L{8.4cm}@{}}
\toprule
Fluxo & Origem & Destino & Tipo \\
\midrule
caminho & \path{plugin\_processor.cpp:145} & \path{pe\_parser.cpp:65} & \texttt{vector<uint8\_t>} \\
arrasto & \path{plugin\_editor.h:66} & \path{plugin\_processor.cpp:144} & \texttt{juce::String} \\
midi e params & \path{plugin\_processor.cpp:471} & \path{note\_tracker.cpp:7} & \texttt{MidiBuffer} e \texttt{float} \\
audio e estado & \path{plugin\_processor.cpp:499} & barramento do host & \texttt{float}, dois canais \\
preset & \path{plugin\_processor.cpp:658} & ficheiro \texttt{.opcoda} & texto XML \\
\bottomrule
\end{tabular}
\end{center}

O ficheiro PE entra pela thread de interface e nunca pela thread de audio. O
que viaja para o motor e' um ponteiro, e o que o motor devolve e' o ponteiro
antigo, para que a interface possa libertar a memoria quando quiser.

\section*{3 Contentores}

\begin{center}
\maxfig{containers}
\end{center}

\begin{center}
\begin{tabular}{@{}L{7.0cm}L{6.4cm}L{11.6cm}@{}}
\toprule
Relacao & Onde & O que significa \\
\midrule
\texttt{opcoda\_core} para \texttt{pe} & \path{CMakeLists.txt:22} & leitura do PE, conversao de bytes, selecao de regiao \\
\texttt{opcoda\_core} para \texttt{entropy} & \path{CMakeLists.txt:27} & entropia de Shannon por janela \\
\texttt{opcoda\_core} para \texttt{dsp} & \path{CMakeLists.txt:33} & motor granular, janelas, filtros, limitador \\
\texttt{opcoda\_core} para \texttt{rt} & \path{CMakeLists.txt:34} & fila sem trava, transporte, notas, guard de alocacao \\
\texttt{opcoda\_core} para \texttt{view} & \path{grain\_trail.h:1} & rastro dos graos, sem JUCE \\
\texttt{opcoda\_plugin} para \texttt{opcoda\_core} & \path{opcoda\_plugin/CMakeLists.txt:82} & o plugin liga o nucleo \\
\texttt{opcoda\_plugin} para \texttt{processor} & \path{opcoda\_plugin/CMakeLists.txt:61} & fronteira com o host \\
\texttt{opcoda\_plugin} para \texttt{editor} & \path{opcoda\_plugin/CMakeLists.txt:63} & interface \\
\texttt{processor} para \texttt{dsp} & \path{plugin\_processor.cpp:356} & fila SPSC, sem alocacao \\
\bottomrule
\end{tabular}
\end{center}

\subsection*{3.1 Onde esta a fronteira de tempo real}

\texttt{PluginProcessor::processBlock} e' o unico ponto onde a thread de
audio entra. A \texttt{PluginProcessor} tem duas filas: \texttt{incoming\_}
leva material novo para o motor, \texttt{returned\_} traz de volta o
ponteiro que o motor ja nao usa. As duas tem capacidade 4 e
\texttt{SpscRing} exige tipo trivial, por isso o que circula e' ponteiro
cru e nao \texttt{shared\_ptr}.

O buffer vive em \texttt{owned\_}, que e' da thread de interface. A interface
liberta-o em \texttt{PluginEditor::refresh}, a 60 Hz, e so depois de o motor
o ter devolvido.

\section*{4 Fluxo de dados}

\subsection*{4.1 Nivel 0}

Da esquerda para a direita: o ficheiro entra, o audio sai. As duas caixas
cinzentas sao as duas threads, e a fila SPSC e' a unica passagem entre elas.

\begin{center}
\maxfig{dataflow-l0}
\end{center}

\begin{center}
\begin{tabular}{@{}L{3.4cm}L{5.4cm}L{7.6cm}L{8.6cm}@{}}
\toprule
Fluxo & Tipo & Origem & Destino \\
\midrule
bytes & \texttt{vector<uint8\_t>} & \path{plugin\_processor.cpp:145} & \path{pe\_parser.cpp:65} \\
bytes & \texttt{vector<uint8\_t>} & \path{plugin\_processor.cpp:145} & \path{byte\_to\_sample.cpp:5} \\
bytes & \texttt{vector<uint8\_t>} & \path{plugin\_processor.cpp:145} & \path{column\_reduction.cpp:54} \\
secoes & \texttt{pe::Section} & \path{pe\_parser.cpp:126} & \path{column\_reduction.cpp:54} \\
amostras & \texttt{vector<float>} & \path{byte\_to\_sample.cpp:14} & \path{plugin\_processor.cpp:352} \\
curva & \texttt{vector<float>} & \path{column\_reduction.cpp:163} & \path{plugin\_processor.cpp:353} \\
fila & \texttt{PublishedMaterial} & \path{plugin\_processor.cpp:356} & \path{granular\_engine.cpp:118} \\
midi & \texttt{MidiBuffer} & \path{plugin\_processor.cpp:471} & \path{note\_tracker.cpp:7} \\
gate & \texttt{bool} & \path{note\_tracker.cpp:48} & \path{granular\_engine.cpp:65} \\
parametros & \texttt{GranularParams} & \path{plugin\_processor.cpp:484} & \path{granular\_engine.cpp:280} \\
audio & \texttt{float}, dois canais & \path{granular\_engine.cpp:478} & \path{plugin\_processor.cpp:499} \\
\bottomrule
\end{tabular}
\end{center}

A curva de entropia viaja com as amostras porque o motor nao a pode calcular.
Medir a entropia de toda a regiao dentro de \texttt{processBlock} seria
trabalho do tamanho da regiao dentro da thread de audio.

\subsection*{4.2 Nivel 1: o ingest}

Do ficheiro ate a fila. Tudo o que acontece aqui corre na thread de
interface, e por isso pode alocar.

\begin{center}
\maxfig{dataflow-l1-ingest}
\end{center}

\begin{center}
\begin{tabular}{@{}L{1.4cm}L{7.0cm}L{16.6cm}@{}}
\toprule
Passo & Codigo & O que move \\
\midrule
1 & \path{plugin\_processor.cpp:144} & o caminho do ficheiro \\
2 & \path{plugin\_processor.cpp:160} & os bytes crus \\
3 & \path{pe\_parser.cpp:65} & a estrutura do PE, ou um erro tipado \\
4 & \path{pe\_parser.cpp:126} & uma \texttt{Section} por seccao \\
5 & \path{shannon\_entropy.cpp:21} & os bits por byte de cada seccao \\
6 & \path{byte\_to\_sample.cpp:14} & as amostras em $[-1, 1]$ \\
7 & \path{column\_reduction.cpp:163} & minimo, maximo, RMS e entropia de cada coluna \\
8 & \path{plugin\_processor.cpp:317} & as amostras e a curva, em dois vectores \\
9 & \path{plugin\_processor.cpp:356} & os dois ponteiros, para a fila \\
10 & \path{plugin\_processor.cpp:167} & o erro, como \texttt{const char*} \\
\bottomrule
\end{tabular}
\end{center}

\subsection*{4.3 Nivel 1: o bloco de audio}

Do host ate ao barramento. Nada disto aloca, nada disto bloqueia, e nada
disto le disco.

\begin{center}
\maxfig{dataflow-l1-audio}
\end{center}

\begin{center}
\begin{tabular}{@{}L{1.4cm}L{7.0cm}L{16.6cm}@{}}
\toprule
Passo & Codigo & O que move \\
\midrule
1 & \path{plugin\_processor.cpp:465} & o buffer e o MIDI do bloco \\
2 & \path{plugin\_processor.cpp:356} & o \texttt{PublishedMaterial} que chegou \\
3 & \path{plugin\_processor.cpp:387} & o ponteiro das amostras e o da curva \\
4 & \path{plugin\_processor.cpp:414} & os eventos MIDI do bloco \\
5 & \path{note\_tracker.cpp:7} & as notas premidas e o pedal \\
6 & \path{transport.cpp:122} & a fracao da regiao que o transporte leu \\
7 & \path{granular\_engine.cpp:65} & o comando de abrir ou fechar o gate \\
8 & \path{granular\_engine.cpp:178} & o grao: posicao, ganho, pan e janela \\
9 & \path{granular\_engine.cpp:392} & a amostra interpolada, ja filtrada \\
10 & \path{granular\_engine.cpp:495} & o nivel do envelope, aplicado a mistura \\
11 & \path{limiter.cpp:62} & o sinal saturado \\
12 & \path{plugin\_processor.cpp:499} & os dois canais no barramento \\
\bottomrule
\end{tabular}
\end{center}

O gate vem antes do limitador de proposito. Fechar a rampa antes de saturar
mantem o tecto de saida independente de quantas notas estao a soar.

\section*{5 Entradas e saidas}

\subsection*{5.1 Entradas}

\begin{longtable}{@{}L{6.4cm}L{8.6cm}L{10.0cm}@{}}
\toprule
Entrada & Sitio & Dado \\
\midrule
\endfirsthead
\toprule
Entrada & Sitio & Dado \\
\midrule
\endhead
dialogo de ficheiro & \path{plugin\_editor.cpp:1430} & caminho, \texttt{juce::String} \\
arrasto de ficheiro & \path{plugin\_editor.h:66} & \texttt{juce::StringArray} \\
estado do host & \path{plugin\_processor.cpp:637} & bloco XML \\
MIDI do host & \path{plugin\_processor.cpp:471} & \texttt{juce::MidiBuffer} \\
parametros do host & \path{plugin\_processor.cpp:83} & \texttt{float}, 28 parametros \\
taxa de amostragem & \path{plugin\_processor.cpp:120} & \texttt{double} \\
fila de entrada & \path{plugin\_processor.cpp:356} & \texttt{PublishedMaterial} \\
fila de devolucao & \path{plugin\_processor.cpp:380} & \texttt{PublishedMaterial} \\
preset do disco & \path{plugin\_processor.cpp:601} & \texttt{juce::XmlElement} \\
ficheiro PE pelo nucleo & \path{pe\_parser.cpp:166} & \texttt{vector<uint8\_t>} \\
\bottomrule
\end{longtable}

\subsection*{5.2 Saidas}

\begin{longtable}{@{}L{6.4cm}L{8.6cm}L{10.0cm}@{}}
\toprule
Saida & Sitio & Dado \\
\midrule
\endfirsthead
\toprule
Saida & Sitio & Dado \\
\midrule
\endhead
audio para o host & \path{plugin\_processor.cpp:499} & \texttt{float}, dois canais \\
estado para o host & \path{plugin\_processor.cpp:595} & bloco XML \\
parametro para o host & \path{plugin\_processor.cpp:457} & \texttt{float} normalizado \\
preset para o disco & \path{plugin\_processor.cpp:658} & texto XML \\
telemetria para a interface & \path{plugin\_processor.cpp:522} & atomicos \\
grafos para a interface & \path{granular\_engine.cpp:492} & \texttt{GrainTelemetry} \\
erro de parser & \path{plugin\_processor.cpp:167} & \texttt{const char*} \\
erro de abertura & \path{plugin\_processor.cpp:147} & \texttt{const char*} \\
erro de fila cheia & \path{plugin\_processor.cpp:359} & \texttt{const char*} \\
erro de fonte em falta & \path{plugin\_processor.cpp:632} & \texttt{const char*} \\
\bottomrule
\end{longtable}

\subsection*{5.3 Erros}

Os erros do Opcoda sao codigos de texto, nao excepcoes.
\texttt{pe::toString} traduz o enumerado \texttt{PeError} para uma cadeia
estavel, e a interface mostra a cadeia tal e qual.

\begin{center}
\begin{tabular}{@{}L{8.0cm}L{17.0cm}@{}}
\toprule
Codigo & Quando \\
\midrule
\texttt{E\_OPEN} & o ficheiro nao abre \\
\texttt{E\_EMPTY} & o ficheiro esta vazio \\
\texttt{E\_TRUNCATED} & a leitura devolveu menos bytes do que o tamanho declarado \\
\texttt{E\_BAD\_MZ} & a assinatura \texttt{MZ} nao esta no inicio \\
\texttt{E\_BAD\_E\_LFANEW} & o offset do cabecalho PE esta fora do ficheiro \\
\texttt{E\_BAD\_PE} & a assinatura \texttt{PE} nao esta no sitio \\
\texttt{E\_BAD\_OPTIONAL\_HEADER} & o cabecalho opcional e' curto demais \\
\texttt{E\_TOO\_MANY\_SECTIONS} & o numero de seccoes e' zero ou passa de 96 \\
\texttt{E\_BAD\_SECTION\_TABLE} & a tabela de seccoes nao cabe no ficheiro \\
\texttt{E\_OOB} & uma seccao declara mais bytes do que o ficheiro tem \\
\texttt{E\_BUSY} & a fila de entrada esta cheia \\
\texttt{E\_SOURCE\_MISSING} & o projecto aponta para um ficheiro que ja nao esta la \\
\bottomrule
\end{tabular}
\end{center}

\subsection*{5.4 Caminho quente}

\begin{center}
\maxfig{sequence-hot-path}
\end{center}

\begin{multicols}{2}
O arrasto e' tratado em \texttt{PluginEditor::filesDropped}, que chama
\texttt{ingestFirst}. O dialogo de ficheiro faz o mesmo por
\texttt{chooseFile}. Os dois caminhos acabam no mesmo
\texttt{PluginProcessor::ingest}, e nao ha caminho paralelo.

A troca de material acontece no inicio do bloco, em
\texttt{drainIncomingQueue}. O motor recebe o ponteiro novo em
\texttt{setSource} e o ponteiro antigo volta pela fila de devolucao. A
interface liberta-o no \texttt{refresh} seguinte, a 60 Hz.
\end{multicols}

\end{document}

